Lucrarea prezintă proiectarea și implementarea platformei \textit{Favigon}, o aplicație web full-stack orientată spre construirea, editarea, partajarea și exportul interfețelor web. Ideea proiectului a pornit de la o observație simplă: în multe fluxuri de lucru actuale, trecerea de la design la cod rămâne fragmentată. Designerul lucrează într-un instrument, dezvoltatorul reconstruiește manual interfața, iar utilizatorul final vede rezultatul abia după mai multe iterații. În cazul proiectelor mici această separare este suportabilă, însă pentru proiecte dinamice devine rapid o sursă de timp pierdut, nealiniere și compromisuri tehnice.

În acest context, lucrarea urmărește realizarea unei platforme care să combine mai multe direcții într-un singur produs: editare vizuală într-un canvas, generare asistată de inteligență artificială, export de cod pentru mai multe framework-uri și o zonă socială în care proiectele pot fi publicate, apreciate și reutilizate. În forma actuală, Favigon se adresează în special designerilor, frontend developerilor, studenților și echipelor mici care vor să transforme rapid o idee de interfață într-un prototip editabil, într-un model intern coerent și, la final, în cod executabil. Din punct de vedere tehnic, aplicația folosește un frontend dezvoltat în Angular 20, un backend REST realizat în ASP.NET Core 9, o bază de date PostgreSQL și o infrastructură containerizată cu Docker Compose. Pentru expunerea aplicației în exterior este folosit și un tunel Cloudflare.

O parte importantă a lucrării este reprezentată de editorul canvas. Acesta nu este tratat ca o simplă interfață de desen, ci ca un subsistem complex, cu stare proprie, istoric de editare, persistare automată, operații de clipboard, gesturi de manipulare și suport pentru mai multe pagini. În jurul lui este construit un model de date orientat pe elemente și pagini, iar pe frontend logica este organizată în servicii dedicate pentru viewport, gesturi, istoric, persistare și generare de cod. Această separare menține codul relativ ușor de extins și evită concentrarea tuturor responsabilităților într-o singură componentă.

Capitolul central al lucrării este dedicat funcționalității de generare asistată de AI și conversiei din design în cod. În locul unei generări directe, greu de controlat, este folosită o soluție în trei faze. Prima fază extrage intenția utilizatorului și definește structura logică a paginii. A doua fază generează o reprezentare intermediară a interfeței, bazată pe un arbore de noduri cu informații de layout, stil, poziționare și metadate. A treia fază aplică stilizarea finală, păstrând structura deja validată. În acest mod, partea de AI poate fi constrânsă mai bine, iar rezultatele sunt mai ușor de verificat și reparat atunci când modelul produce ieșiri imperfecte.

Conversia în cod folosește aceeași reprezentare intermediară și o transformă în fișiere pentru HTML, React și Angular. Astfel, editorul și motorul de export nu sunt legate rigid de un singur framework. Din punct de vedere practic, această decizie a fost una dintre cele mai importante din proiect, deoarece a separat clar problema construirii designului de problema randării lui într-o tehnologie concretă. În plus, platforma suportă și export multi-page, precum și generare responsive prin diferențe între breakpoints.

Lucrarea tratează și componentele conexe care fac aplicația utilizabilă ca produs, nu doar ca demo tehnic: autentificare cu JWT în cookie-uri, OAuth prin furnizori externi, resetare parolă, autentificare în doi pași, gestiunea profilului, proiecte publice și private, explorare de proiecte și funcții sociale precum like, star, follow sau fork. În plus, sunt prezentate măsurile de securitate aplicate în backend, cum ar fi rate limiting-ul, limitarea dimensiunii cererilor, middleware-ul de tratare a excepțiilor și politicile de securitate HTTP.

Pentru validare, proiectul include atât testare automată, cât și scenarii funcționale. La momentul redactării lucrării, suita de teste backend conține 108 teste care verifică servicii, controllere, middleware și motorul de conversie. Partea de frontend este testată mai puțin automatizat, motiv pentru care în lucrare sunt prezentate și validarea prin scenarii operaționale, observații directe și demonstrarea fluxurilor principale.

Rezultatul final este o platformă care nu încearcă să înlocuiască instrumentele consacrate din ecosistem, ci să ofere un flux integrat pentru proiectare rapidă, experimentare, export și publicare. Aportul personal se reflectă mai ales în integrarea coerentă a acestor module, în modelul intermediar folosit pentru design și cod, în pipeline-ul AI în trei faze și în modul în care editorul canvas a fost separat pe servicii specializate. În ansamblu, proiectul a permis combinarea unor teme de arhitectură software, dezvoltare web, securitate, UX și inteligență artificială aplicată într-un produs unitar.
